% concatenated from main.tex
\documentclass[11pt]{amsart}
\usepackage{epsfig,color}
\usepackage{blindtext}
\usepackage{graphicx}
\usepackage{enumitem}
\usepackage{url}
\usepackage{amssymb}
\usepackage{graphicx,import}
\usepackage{comment}
\usepackage{esint}
\usepackage{xypic}
\usepackage{amsthm}

%% page size
%\headheight=6.15pt \textheight=8in \textwidth=6.5in
%\oddsidemargin=0in \evensidemargin=0in \topmargin=0in
\usepackage[margin = 1.4in] {geometry}

\usepackage{hyperref}

%------------------------------------------------------------------------------------%
\newtheorem*{restate}{Theorem}
\newtheorem{thm}{Theorem}[section]
\newtheorem{cor}[thm]{Corollary}
\newtheorem{lem}[thm]{Lemma}
\newtheorem{prop}[thm]{Proposition}
\newtheorem{defn}[thm]{Definition}
\newtheorem{rem}[thm]{\bf{Remark}}
\newtheorem{alg}{\bf{Algorithm}}
\newtheorem*{Proof}{Proof}
\newtheorem{ex}[thm]{Example}
\numberwithin{equation}{section}
\def\pn{\par\noindent}
\def\cen{\centerline}
%------------------------------------------------------------------------------------%
% theorem environments
\setcounter{section}{0}
\theoremstyle{theorem}

\newtheorem{theorem}{Theorem}[section]
\newtheorem{conjecture}[theorem]{Conjecture}
\newtheorem{proposition}[theorem]{Proposition}
\newtheorem{lemma}[theorem]{Lemma}
\newtheorem{corollary}[theorem]{Corollary}
\newtheorem{question}[theorem]{Question}
\newtheorem{fact}[theorem]{Fact}

\theoremstyle{definition}
\newtheorem{example}[theorem]{Example}
\newtheorem{definition}[theorem]{Definition}
\newtheorem*{remark*}{Remark}
\newtheorem*{acknowledgements}{Acknowledgements}
\newtheorem{remark}[theorem]{Remark}


\numberwithin{equation}{section}


% commands
\newcommand{\cA}{\mathcal{A}}
\renewcommand{\cH}{\mathcal{H}}
\newcommand{\R}{\mathbb{R}}
\newcommand{\Q}{\mathbb{Q}}
\newcommand{\N}{\mathbb{N}}
\newcommand{\Z}{\mathbb{Z}}

\DeclareMathOperator{\Ric}{Ric}
\DeclareMathOperator{\Lip}{Lip}
\DeclareMathOperator{\diam}{diam}
\newcommand{\eps}{\varepsilon}
\DeclareMathOperator{\Image}{image}
\DeclareMathOperator{\codim}{codim}
\DeclareMathOperator{\inj}{inj}
\DeclareMathOperator{\vol}{Vol} 


\author{Elham Matinpour}
\address{}
\email{}


\subjclass[2020]{Primary 53A10; Secondary 49Q05, 53C42}
\keywords{free-boundary minimal surfaces, Plateau singularities, tetrahedral cone, conformal minimal immersions, rigidity}


\title{Free Boundary Plateau Model Cones in $\mathbb{B}^n$ are Rigid under Conformal Minimal Immersions}


\begin{document}


\maketitle

% ABSTRACT
\begin{abstract}
The classical theorem of Nitsche~\cite{nitsche1985stationary} asserts that every free-boundary minimal disk in the unit ball \(\mathbb{B}^3\) is an equatorial plane disk. Fraser and Schoen~\cite{fraser2015uniqueness} extended this rigidity result to all dimensions and to spaces of constant sectional curvature, showing in particular that such disks are totally geodesic. In a previous work~\cite{matinpour2025cone}, the author established an analogue for singular \(Y\)-type surfaces: any conformal minimal immersion of the flat \(Y\)-cone into \(\mathbb{B}^n\) that meets the boundary sphere orthogonally must coincide with the flat \(Y\)-cone itself (up to orthogonal transformation).

In the present paper we push this program one step further by treating the more intricate tetrahedral singularity. We prove that every conformal and minimal immersion of the flat \(T\)-cone into the unit ball \(\mathbb{B}^n\), subject to the free-boundary orthogonality condition, is necessarily congruent to the flat \(T\)-cone. Combining this with the smooth-disk result of Nitsche--Fraser--Schoen and the \(Y\)-cone result of~\cite{matinpour2025cone}, we obtain a unified uniqueness theorem: any free-boundary minimal Plateau surface in \(\mathbb{B}^n\) that is conformal to one of the reference model domains (a disk in the plane, the \(Y\)-cone, or the \(T\)-cone) must itself be congruent to that model.
\end{abstract}


\section{Introduction}
This paper continues the line of investigation initiated in~\cite{matinpour2025cone}, where the uniqueness of the flat $Y$-cone was established among all free-boundary minimal $Y$-surfaces that are conformal to the model $Y$-cone. Here we address the next natural singular model---the tetrahedral \(T\)-cone---and thereby complete the picture for all admissible Plateau-type singularities under the conformality assumption.\\
In the smooth (non-singular) setting the foundational rigidity result is due to Nitsche~\cite{nitsche1985stationary}, who proved that any free-boundary minimal disk in \(\mathbb{B}^3\) is an equatorial flat disk. Fraser and Schoen~\cite{fraser2015uniqueness} later generalized this statement to arbitrary dimension and to ambient manifolds of constant curvature showing that such disks are totally geodesic (and, under an additional parallel mean-curvature assumption, totally umbilic in a three-dimensional constant-curvature submanifold). These theorems provide the regular building block of our program.\\
For Plateau surfaces the situation is substantially more delicate because of the presence of the classical \(Y\)- and \(T\)-junction singularities. The treatment of singular cases is settled by exploiting the conformal structure on each sector together with the minimality configuration along the singular curves and the free boundary that forces planarity of every face. 

Our main results are the following two theorems.

\begin{theorem}\label{thm:uniq.TC}
Let \(TC\) denote the flat \(T\)-cone. Suppose
\[
u : TC \to \mathbb{B}^n
\]
is a conformal minimal immersion such that \(u(TC)\) meets \(\partial \mathbb{B}^n\) orthogonally. Then \(u(TC)\) coincides with the flat \(T\)-cone up to an orthogonal transformation, i.e.,
\[
u(TC) = A(TC) \quad \text{for some } A \in \mathrm{O}(n).
\]
\end{theorem}
Combining Theorem~\ref{thm:uniq.TC} with the classical rigidity theorem of Nitsche--Fraser--Schoen for free-boundary minimal disks and the uniqueness result for the flat \(Y\)-cone established in~\cite{matinpour2025cone}, we obtain a unified rigidity theorem for all classical two-dimensional Plateau model minimal surfaces.

\begin{theorem}\label{thm:uniq.PlateauSurfaces}
Let \(\Sigma\) be a free-boundary minimal Plateau surface in \(\mathbb{B}^n\) arising as the image of a conformal minimal immersion of one of the classical model domains: a planar disk, the flat \(Y\)-cone, or the flat \(T\)-cone. Then \(\Sigma\) is congruent, via an orthogonal transformation of \(\mathbb{R}^n\), to the corresponding flat model.
\end{theorem}

Thus, under the conformality assumption, the standard flat disk, the flat \(Y\)-cone, and the flat \(T\)-cone are rigid among all free-boundary minimal Plateau surfaces modeled on the corresponding singularity type. In this way, the present work completes a unified rigidity theory for the classical smooth, \(Y\)-type, and tetrahedral Plateau configurations.\\
For the precise definitions of Plateau surfaces, the model cones, the free-boundary condition, and the notion of conformality used throughout the paper, we refer the reader to the Preliminaries (Section~\ref{Sec1-Preliminaries}).

The remainder of the paper is organized as follows. In Section~\ref{sec:uniq-tsg-net}, we show that any stationary geodesic network on $\mathbb{S}^n$ with tetrahedral combinatorics coincides with the regular tetrahedral network. In Section~\ref{sec:uniq. FBM TS with NT}, we establish that free-boundary minimal $T$-surfaces spanning a regular tetrahedral network on $\mathbb{S}^n$ are uniquely given by the regular tetrahedral cone $TC$. Section~\ref{sec:uniq. Tcone} completes the proof of Theorem~\ref{thm:uniq.TC} by analyzing the Hopf differential on the faces, showing that the tetrahedral network is composed of geodesic arcs, and then applying the rigidity result from Section~\ref{sec:uniq. FBM TS with NT}. Finally, the global result (Theorem~\ref{thm:uniq.PlateauSurfaces}) follows by invoking known results for the regular and $Y$-cases.

\subsection*{Acknowledgment} I would like to thank Professor Jaigyoung Choe for suggesting the problem.




\section{Preliminaries}\label{Sec1-Preliminaries}

We consider two-dimensional rectifiable currents (more precisely, integral currents) in $\mathbb{R}^n$. A rectifiable current $\Sigma$ is said to be a \emph{Plateau surface} if its support satisfies the following local structural conditions.

\begin{itemize}
    \item The support of $\Sigma$ is locally $C^{1,\alpha}$-diffeomorphic (for some $\alpha \in (0,1)$) to one of the model cones: a plane $P$, a half-plane $H$, a $Y$-cone $Y$, or a $T$-cone $T$. Here:
    \begin{itemize}
        \item $P$ is a $2$-dimensional affine plane in $\mathbb{R}^n$ (the local model for boundary points and free-boundary arcs),
        \item $H$ is a $2$-dimensional half-plane (the local model for boundary points and free-boundary arcs),
        \item $Y$ is the union of three half-planes meeting along a common line at $120^\circ$ angles (the local model for $Y$-junctions),
        \item $T$ is the cone over the $1$-skeleton of a regular tetrahedron centered at the origin (the local model for $T$-points).
    \end{itemize}
    More precisely, for every point $p \in \mathrm{spt}(\Sigma)$ there exist a neighborhood $N(p) \subset \mathrm{spt}(\Sigma)$ and a $C^{1,\alpha}$ diffeomorphism $u : N(p) \to u(N(p))$ mapping $N(p)$ onto one of the four model cones.
    
    \item At each point $p \in \mathrm{spt}(\Sigma)$, the tangent cone $\mathrm{Tan}(\Sigma, p)$, in the varifold or current sense, coincides (up to orthogonal transformation) with one of the four model cones $P$, $H$, $Y$, or $T$.
\end{itemize}

A Plateau surface $\Sigma$ is called \emph{minimal} if it is area-stationary, i.e., a stationary point for the area functional, or equivalently, it has a weakly vanishing mean curvature at regular points.
\medskip

We further assume that $\Sigma$ is \emph{two-sided}, meaning that it separates in $\mathbb{R}^n$ locally: around every interior point there is a small ball in which $\Sigma$ divides the ball into exactly two connected components.\\
More globally, $\Sigma$ defines a cell decomposition of an open set $U \subset \mathbb{R}^n$ if there is a finite or countable family of open connected sets $\{U^i\}$ (the \emph{cells}) such that $U \setminus \mathrm{spt}(\Sigma) = \bigcup_i U^i$ (disjoint union), $\partial \Sigma \subset \partial U$, and for every $p \in \mathrm{spt}(\Sigma) \cap U$ there exists $r>0$ such that for all $0 < r' < r$ and every cell $U^i$, $B_{r'}(p) \cap U^i$ is connected (possibly empty). Here, $B_r(p)$ is the ball with radius $r$ and centered at $p$ in $\mathbb{R}^n$.
\medskip

In this paper we focus on \emph{$T$-surfaces}: minimal Plateau surfaces that have at least one $T$-point.

Points are classified as follows:
\begin{itemize}
    \item $p$ is an \emph{interior point} if $\mathrm{Tan}(\Sigma, p) \cong P$ (regular point),
    \item $p$ is a \emph{boundary point} if $\mathrm{Tan}(\Sigma, p) \cong H$,
    \item $p$ is a \emph{$Y$-point} if $\mathrm{Tan}(\Sigma, p) \cong Y$,
    \item $p$ is a \emph{$T$-point} if $\mathrm{Tan}(\Sigma, p) \cong T$.
\end{itemize}

We denote:
\begin{align*}
    \mathrm{reg}(\Sigma) &:= \{p \in \Sigma : p \text{ is interior}\}, \\
    \mathcal{R} &:= \text{connected components of } \mathrm{reg}(\Sigma) \quad (\text{the regular faces}), \\
    \partial \Sigma &:= \{p \in \Sigma : p \text{ is a boundary point}\}, \\
    \mathcal{B} &:= \text{connected components of } \partial \Sigma \quad (\text{boundary arcs}), \\
    \Sigma_Y &:= \{p \in \Sigma : p \text{ is a } Y\text{-point}\}, \\
    \mathcal{Y} &:= \text{connected components of } \Sigma_Y \quad (\text{$Y$-junctions}), \\
    \Sigma_T &:= \{p \in \Sigma : p \text{ is a } T\text{-point}\}.
\end{align*}
A \emph{$T$-surface with a single $T$-point} has exactly one point in $\Sigma_T$, four $Y$-junctions emanating from the $T$-point, and six regular faces. We work with compact Plateau surfaces with boundary, naturally modeled as an integral current.\\
Define the quarter-disk
\[
\hat{D} = \{ (x,y) \in \mathbb{R}^2 : x^2 + y^2 \leq 1,\ x \geq 0,\ y \geq 0 \},
\]
with boundary parts
$$\gamma_1 = \{ (0,y) : 0 \leq y \leq 1 \},\ \
\gamma_2 = \{ (x,0) : 0 \leq x \leq 1 \},\ \
\sigma = \{ (x,y) : x^2 + y^2 = 1,\ x \geq 0,\ y \geq 0 \}.$$
Let $P$ be a regular tetrahedron in $\mathbb{R}^n$ centered at the origin with vertices $v_1,\dots,v_4 \in \mathbb{S}^{n-1}$. Let $E(P)$ be its $1$-skeleton (six edges). The \emph{flat $T$-cone} is
\[
TC = \{ t x : x \in E(P),\ 0 \leq t \leq 1 \},
\]
the union of six flat sectors from the origin over the edges of $P$. The four radial segments
\[
\Gamma_i = \{ t v_i : 0 \leq t \leq 1 \}, \quad i=1,\dots,4,
\]
are the \emph{$Y$-junctions} of $TC$, along each of which three faces meet at $120^\circ$.\\
We write 
$$TC = \bigcup_{j=1}^6 F_j,$$ where each $F_j$ is a flat sector diffeomorphic to $\hat{D}$, with $\partial F_j = \Gamma_{j_1} \cup \Gamma_{j_2} \cup \partial' F_j$ so that $\Gamma_{j_1},\Gamma_{j_2}$ are junctions and $\partial' F_j$ is a boundary arc on $\mathbb{S}^{n-1}$.

\begin{definition}
A map $\varphi = (\varphi_1,\dots,\varphi_6) : TC \to \mathbb{R}^n$ is a \emph{minimal immersion} if:
\begin{enumerate}
    \item[(i)] each $\varphi_j : F_j \to \mathbb{R}^n$ is a smooth minimal immersion,
    \item[(ii)] along each junction $\Gamma_k$ the three adjacent faces have the same image $\tilde{\Gamma}_k = \varphi(\Gamma_k)$,
    \item[(iii)] the outward unit conormals $\tau_1,\tau_2,\tau_3$ of the three faces along $\tilde{\Gamma}_k$ satisfy $\tau_1 + \tau_2 + \tau_3 = 0$,
    \item[(iv)] at the vertex the local configuration is that of the $T$-cone.
\end{enumerate}
The image $\Sigma := \varphi(TC) = \bigcup_{j=1}^6 \Sigma_j$ (with junctions $\tilde{\Gamma}_k$) is then a \emph{minimal $T$-surface}.
\end{definition}

\begin{definition}
A minimal $T$-surface $\Sigma = \bigcup_{j=1}^6 \Sigma_j$ in the unit ball $\mathbb{B}^n$ is \emph{free-boundary} if:
\begin{enumerate}
    \item each $\Sigma_j$ is minimally immersed in $\mathbb{B}^n$,
    \item $\partial \Sigma \subset \mathbb{S}^{n-1}$,
    \item $\Sigma$ meets $\mathbb{S}^{n-1}$ orthogonally,
    \item along each junction the conormal balance $\tau_1 + \tau_2 + \tau_3 = 0$ holds,
    \item the vertex configuration is tetrahedral.
\end{enumerate}
\end{definition}

We equip $TC$ with the flat Euclidean metric induced from $\mathbb{R}^n$, inducing compatible metrics $g_j$ on each $F_j$ that agree along junctions.\\
The $T$-surface $\Sigma$ can be equipped with an abstract metric,
not necessarily induced by a Euclidean metric.
Let $g_j$ be a Riemannian metric on $\Sigma_j$.
\medskip

We say that \emph{$(g_1,\ldots,g_6)$ defines a metric on $\Sigma$}
if the induced metrics agree along each junction curve.
More precisely, if three faces $\Sigma_{i},\Sigma_{j},\Sigma_{k}$
meet along $\Gamma_\ell$, then
\[
g_i|_{\Gamma_\ell} = g_j|_{\Gamma_\ell} = g_k|_{\Gamma_\ell}.
\]
Equivalently, all three faces induce the same metric
\[
g_{\Gamma_\ell}
\]
on the junction curve $\Gamma_\ell$.\\
Given such a metric $g$ on $\Sigma$, we define the unit outer conormal
vector field $\tau_j$ to $\Sigma_j$ along each junction curve.\\
Let $\eta$ be the unit tangent vector field along a junction curve
$\Gamma_\ell$.
If the three faces $\Sigma_i,\Sigma_j,\Sigma_k$ meet along $\Gamma_\ell$,
their geodesic curvatures are defined by
\[
\kappa_i = g_i(\nabla_{\eta}^{\Sigma_i} \eta ,\tau_i), \quad
\kappa_j = g_j(\nabla_{\eta}^{\Sigma_j} \eta ,\tau_j), \quad
\kappa_k = g_k(\nabla_{\eta}^{\Sigma_k} \eta ,\tau_k).
\]

\begin{definition}
We say that the metric $g$ on $\Sigma$ is \emph{compatible with the
$T$–structure} if along each junction curve $\Gamma_\ell$ the geodesic
curvatures satisfy
\[
\kappa_i + \kappa_j + \kappa_k = 0 .
\]
\end{definition}

\begin{definition}
A map
\[
u=(u_1,\ldots,u_6):TC \to \mathbb{R}^n
\]
is called \emph{conformal} if each
\[
u_j : F_j \to \mathbb{R}^n
\]
is conformal with respect to $g_j$ and the induced metric on $u_j(F_j)$, and the maps agree on junctions: 
\[
u_i(p)=u_j(p)=u_k(p),
\]
for $p$ on a junction shared by $F_i,F_j,F_k$.
\end{definition}


\section{Rigidity of Tetrahedral Stationary Geodesic Networks}
\label{sec:uniq-tsg-net}

In this section we show that a stationary geodesic network on the sphere with
tetrahedral combinatorics must coincide with the regular tetrahedral
network.

\begin{definition}
Let $\mathcal{N}$ be a finite union of arcs on the unit sphere
$\mathbb S^{n-1}$ meeting at vertices.
We say that $\mathcal{N}$ is a \emph{stationary geodesic network}
if the following conditions hold:

\begin{enumerate}
\item Each edge of $\mathcal{N}$ is a geodesic arc of $\mathbb S^{n-1}$.

\item At each vertex exactly three edges meet.

\item (120$^\circ$ equilibrium condition)
If $v$ is a vertex and $\tau_1,\tau_2,\tau_3$ are the unit tangent
vectors of the three edges at $v$, then
\[
\tau_1+\tau_2+\tau_3=0 .
\]
\end{enumerate}
\end{definition}

\begin{definition}
A stationary geodesic network $\mathcal N\subset \mathbb S^{n-1}$ is said
to have \emph{tetrahedral combinatorics} if

\begin{itemize}
\item it has four vertices,
\item six edges,
\item each pair of vertices is connected by an edge.
\end{itemize}
\end{definition}

\begin{definition}
    Let $\mathcal{NT}\subset \mathbb S^{n-1}$ be a tetrahedral network, we say that $\mathcal{NT}$ is \emph{the regular tetrahedral network} if 

    \begin{itemize}
        \item $\mathcal{NT}\subset \mathbb{S}^2 \subset \mathbb S^{n-1},$
        \item $\mathcal{NT}=TC\cap \mathbb{S}^2,$
    where $TC$ is the $T$-cone defined in Section 2.
         \end{itemize}
       These two conditions are equivalent to for any two vertices
       $$
       v_i\cdot v_j=-\frac{1}{3},\ \ \  (i\neq j).
       $$
       In particular, $\mathcal{NT}$ is a stationary geodesic network which has a tetrahedral combinatorics.  
\end{definition}
 
\begin{lemma}\label{lem: uniq. of RT-network}
Let $\mathcal{N} \subset \mathbb{S}^{n-1}$ be a stationary geodesic network with tetrahedral combinatorics, i.e., it has exactly four vertices, six geodesic edges, and every pair of vertices is connected by an edge. Then $\mathcal{N}$ is isometric to the regular tetrahedral network, that is, there exists an orthogonal transformation $A \in \mathrm{O}(n)$ such that
$
\mathcal{N} = A \bigl( \mathcal{NT} \bigr)$.\\
In particular, all edges of $\mathcal{N}$ have the same spherical length $\arccos(-1/3)$.
\end{lemma}
This result is classical in the context of Steiner networks and 
spherical geometry; we include a proof for completeness.
\begin{proof}
Let $V = \{v_1, v_2, v_3, v_4\} \subset \mathbb{S}^{n-1}$ be the four vertices of $\mathcal{N}$ and let $E$ be the set of six geodesic arcs connecting every pair of vertices.
\medskip

\textit{First,} we show that the network lies in a great 2-sphere. \\ 
Fix any vertex $v \in V$. By definition of a stationary geodesic network, exactly three edges emanate from $v$ and their unit tangent vectors $\tau_1, \tau_2, \tau_3 \in T_v \mathbb{S}^{n-1} \simeq v^\perp$ satisfy
\[
\tau_1 + \tau_2 + \tau_3 = 0.
\]
Thus $\tau_1, \tau_2, \tau_3$ are linearly dependent and span a subspace of dimension at most $2$. Let $T_v \subset v^\perp$ be this $2$-dimensional subspace. Each of the three neighboring vertices $w_i$ ($i=1,2,3$) satisfies
\[
w_i = \cos \theta_i \, v + \sin \theta_i \, u_i, \qquad u_i \in \mathbb{S}^{n-1} \cap v^\perp,
\]
where $u_i$ is parallel to the initial tangent vector of the geodesic from $v$ to $w_i$. Hence $u_i \in T_v$, up to scaling, so $w_i \in \operatorname{span}\{v\} \oplus T_v$, a $3$-dimensional subspace of $\mathbb{R}^n$. Since this holds for every vertex and the graph is connected, all four vertices $v_1,v_2,v_3,v_4$ lie in the same $3$-dimensional subspace $W \subset \mathbb{R}^n$. Therefore
\[
\mathcal{N} \subset W \cap \mathbb{S}^{n-1} \cong \mathbb{S}^2,
\]
It is therefore enough to prove the statement for $n=3$.
\medskip

\textit{Second,} we show that the complement consists of four equiangular spherical triangles. \\ 
The six geodesic arcs on $\mathbb{S}^2$ form the $1$-skeleton of the complete graph $K_4$. Since edges are shortest geodesics, the graph is embedded without crossings. By Euler's formula,
\[
V - E + F = 2 \quad \Rightarrow \quad 4 - 6 + F = 2 \quad \Rightarrow \quad F = 4.
\]
Thus the network bounds exactly four spherical regions, each a triangle. At every vertex, the three incident edges meet at $120^\circ$ by the equilibrium condition, so each spherical triangle is \emph{equiangular} with all angles equal to $120^\circ$.
\medskip


\textit{Third,} we show that each spherical triangle is equilateral. \\ 
Consider one of the four spherical triangles with angles $A = B = C = 120^\circ$. By the spherical law of sines,
\[
\frac{\sin a}{\sin A} = \frac{\sin b}{\sin B} = \frac{\sin c}{\sin C}.
\]
Since $A = B = C$, it follows that $\sin a = \sin b = \sin c$. As each side length lies in $(0,\pi)$, where $\sin$ is injective, we conclude that
\[
a = b = c.
\]
Thus each spherical triangle is equilateral. Since every edge of $\mathcal{N}$ is a side of one of these triangles, all six edges have the same length $\sigma \in (0,\pi)$.
\medskip


\textit{Fourth,} we determine the value of the common side length and identify the configuration. \\ 
Applying the spherical law of cosines to an equilateral spherical triangle with angles $120^\circ$, we obtain
\[
\cos \sigma = \cos^2 \sigma + \sin^2 \sigma \cos(120^\circ),
\]
which simplifies to $\cos \sigma = -\tfrac{1}{3}$ (excluding the degenerate case). Hence $\sigma = \arccos(-1/3)$, and the four vertices form a regular tetrahedron.
\medskip


\textit{At last,} the degenerate cases ($\sigma = 0$ or $\sigma = \pi$) are excluded because they would violate either the finite union of arcs condition, the three-edges-per-vertex condition, or the proper embedding requirement. Thus, the proof is complete.
\end{proof}



\section{Rigidity of Free Boundary Minimal $T$-surfaces spanning  the Regular Tetrahedral  Network}\label{sec:uniq. FBM TS with NT}

\begin{theorem}\label{Thm:uniq. FBMT with NT}
Let \(\Sigma \subset \mathbb{B}^n\) be a free-boundary minimal $2$-dimensional $T$-surface with exactly one $T$-point, four $Y$-junction edges (along each of which three faces meet at $120^\circ$), and six faces. Assume that
\begin{itemize}
\item \(\partial \Sigma \cap \mathbb{S}^{n-1}\) is a stationary geodesic network with tetrahedral combinatorics,
\item \(\Sigma\) meets \(\mathbb{S}^{n-1}\) orthogonally.
\end{itemize}
Then there exists \(A \in \mathrm{O}(n)\) such that \(\Sigma = A(TC)\), where \(TC\) is the $T$-cone over the regular tetrahedron centered at the origin.
\end{theorem}

\begin{proof}
By the uniqueness of stationary geodesic networks with tetrahedral combinatorics (Lemma \ref{lem: uniq. of RT-network}), after applying an orthogonal transformation, we may assume that \(\partial \Sigma \cap \mathbb{S}^{n-1} = \mathcal{NT}\), the regular tetrahedral network lying in a $3$-dimensional subspace $W \subset \mathbb{R}^n$.
\medskip

Let \(p \in \Sigma\) be the unique $T$-point and let \(F\) be any one of the six faces of \(\Sigma\). Then \(\partial F = \sigma \cup \Gamma_1 \cup \Gamma_2\), where \(\sigma \subset \mathcal{NT}\) is a geodesic arc on \(\mathbb{S}^{n-1}\) and \(\Gamma_1, \Gamma_2\) are $Y$-junction curves emanating from \(p\).\\
Let \(\Pi\) be the $2$-plane in \(\mathbb{R}^n\) containing the origin and the great-circle arc \(\sigma \subset \mathcal{NT}\). Since \(\sigma\) is a geodesic arc on \(\mathbb{S}^{n-1}\), \(\Pi\) is invariant under reflection across the unit sphere. By the reflection principle for free-boundary minimal surfaces meeting the sphere orthogonally, Choe \cite{choe2025free}, we may reflect the face \(F\) across \(\mathbb{S}^{n-1}\) along \(\sigma\) to obtain a minimal surface \(F^*\) such that the union
\[
\widetilde{F} := F \cup F^*
\]
is a smooth, in fact real-analytic, minimal surface in a neighborhood of \(\sigma\), with \(\sigma\) now lying in the interior of \(\widetilde{F}\).
\medskip

Because \(F\) meets \(\mathbb{S}^{n-1}\) orthogonally along the boundary arc \(\sigma\), the free-boundary condition implies that the outward conormal of \(\partial F\) in \(F\) coincides with the outward unit normal of the ball, which is the position vector \(q\). In other words,
\[
q \in T_q F \quad \text{for every } q \in \sigma.
\]
Moreover, the tangent vector \(\tau\) to the geodesic arc \(\sigma\) at \(q\) also belongs to \(T_q F\) (since \(\sigma \subset \partial F\)).\\
On other hand,  \(T_q \Pi = \Pi\), that is because \(\Pi\) passes through the origin, and \(\Pi\) is precisely the $2$-plane spanned by the two vectors \(q\) and \(\tau\) (note that \(\tau \perp q\) since \(\tau \in T_q \mathbb{S}^{n-1}\)). Therefore the two vectors \(q\) and \(\tau\) span exactly \(T_q \Pi\), so
\[
T_q F = T_q \Pi \quad \text{for all } q \in \sigma.
\]
The same holds for the reflected piece: \(T_q F^* = T_q \Pi\) for all \(q \in \sigma\) (because the reflection across \(\mathbb{S}^{n-1}\) preserves \(\Pi\) and respects the orthogonal condition). Therefore \(\widetilde{F}\) and the totally geodesic plane \(\Pi\) share the same tangent plane along the entire curve \(\sigma\).
\medskip

Now consider the squared distance function to the plane \(\Pi\),
\[
\rho(x) := \mathrm{dist}^2(x, \Pi),
\]
restricted to \(\widetilde{F}\). Since \(\widetilde{F}\) is minimal and \(\Pi\) is a linear subspace, the function \(\rho\) is subharmonic on \(\widetilde{F}\), this follows from the fact that the Hessian of \(\rho\) along minimal surfaces satisfies the required sign condition; see e.g.\ Colding--Minicozzi \cite{colding2011course} or standard computations for distance functions to affine subspaces.\\
Along the curve \(\sigma\) %(now an interior curve of \(\widetilde{F}\)),     
we have
\[
\rho \equiv 0 \quad \text{and} \quad \nabla \rho \equiv 0 \quad %\text{(since tangent planes agree)}.
\]
Thus \(\rho\) vanishes to first order along an open segment of the interior of \(\widetilde{F}\).\\
Since \(\widetilde{F}\) is real-analytic and $\rho$ is real-analytic, and since $\rho$ vanishes to first order along the interior curve $\sigma$, it follows by unique continuation that $\rho \equiv 0$ in a neighborhood of $\sigma$.
In fact, since the subharmonic \(\rho \geq 0\) everywhere, if \(\rho\) is not identically zero on the connected component of \(\widetilde{F}\) containing \(\sigma\), then by the unique continuation/strong maximum principle for subharmonic functions, %(or the Hopf boundary point lemma applied at points of \(\sigma\)),                 
\(\rho\) cannot attain its minimum $0$ along an interior set unless it is constantly zero in a neighborhood. But \(\rho \equiv 0\) along \(\sigma\), so we conclude that
\[
\rho \equiv 0 \quad \text{in a neighborhood of } \sigma \text{ inside } \widetilde{F}.
\]
In other words, \(\widetilde{F}\) lies in \(\Pi\) near \(\sigma\).\\
Since \(\widetilde{F}\) is a real-analytic minimal surface in the interior and it coincides with \(\Pi\) on an open set, by the unique continuation principle for real-analytic functions %(or identity theorem for minimal surfaces), 
we have
\[
\widetilde{F} \subset \Pi
\]
in the connected component containing \(\sigma\). As \(F\) is connected and contained in one side of the reflection, it follows that
\[
F \subset \Pi.
\]
This holds for every face of \(\Sigma\). Therefore each of the six faces lies in a $2$-plane through the origin. Given that the $Y$-junction curves are minimal and connect the unique $T$-point to the boundary vertices while satisfying the $120^\circ$ condition in each adjacent face, and now knowing the faces are planar, the only possibility is that each $Y$-junction is a straight line segment from the $T$-point to the corresponding boundary vertex.\\ %(by the convexity of geodesics in flat planes and the rigidity of planar triple junctions at $120^\circ$).
Finally, Since the surface $\Sigma$ is a free boundary surface in $\mathbb{B}^n$, the $Y$-junctions must meet the sphere $\mathbb{S}^{n-1}$ orthogonally, which implies that these straight line segments must pass through the origin. Thus, the single $T$-point, as the vertex of the four edges formed with the straight $Y$-junctions, must be the origin. Therefore, each junction is radial and \(\Sigma\) is precisely the flat $T$-cone over the regular tetrahedral network \(\mathcal{NT}\), as desired. 
\end{proof}


\section{Rigidity of the Free-Boundary Minimal Tetrahedral Cone}\label{sec:uniq. Tcone}

\begin{theorem}\label{thm:Tcone}
Let $TC = \bigcup_{j=1}^{6} F_j$ be the flat $T$--cone. Suppose 
$u : TC \to \mathbb{B}^n$ is a conformal minimal immersion such that 
$u(TC)$ meets $\partial \mathbb{B}^n$ orthogonally. Then \[
u(TC) = A(TC) \quad \text{for some } A \in \mathrm{O}(n).
\]
\end{theorem}


\begin{proof}
    The argument follows the same strategy as the uniqueness proof for the flat $Y$-cone in~\cite{matinpour2025cone}, adapted to the tetrahedral combinatorics of the $T$-cone. We work on each quarter-disk sector of the domain and construct a holomorphic function on that sector whose boundary behavior, combined with the junction conditions, forces the relevant second derivatives to vanish along the free-boundary arcs. This implies that those arcs are geodesic on~$\mathbb{S}^{n-1}$, reducing the problem to the rigidity statement of Theorem~\ref{Thm:uniq. FBMT with NT}.
\medskip

Let $TC = \bigcup_{j=1}^{6} F_j$ be the flat $T$-cone consisting of six planar faces $F_j$, each parameterized by the quarter-disk $\hat{D}\subset\mathbb{C}$ (with polar coordinates $(r,\theta)$) meeting along four junction curves $\Gamma_j$ in the tetrahedral configuration. Suppose
\[
u = (u_1,\dots,u_6) : TC \to \mathbb{B}^n
\]
is a conformal minimal immersion such that $u(TC)$ meets $\partial\mathbb{B}^n$ orthogonally. Write
\[
\Sigma = u(TC) = \bigcup_{j=1}^{6} \Sigma_j, \quad \Sigma_j = u_j(\hat{D}),
\]
where each $\Sigma_j$ is a smooth minimal surface. The boundary $\partial\Sigma$ lies on the unit sphere $\mathbb{S}^{n-1}$ and meets it orthogonally.
\medskip

We work throughout in the Euclidean inner product on $\mathbb{R}^n$, extended $\mathbb{C}$-bilinearly to $\mathbb{C}^n$. On each quarter-disk $\hat{D}$, the map $u_j$ is harmonic and conformal:
\[
(u_j)_{z\bar{z}}=0, \qquad (u_j)_z\cdot(u_j)_z=0.
\]
In polar coordinates the Laplacian reads
\[
\Delta = \partial_{rr} + \frac{1}{r}\partial_r + \frac{1}{r^2}\partial_{\theta\theta},
\]
so harmonicity is equivalent to
\[
(u_j)_{rr} + \frac{1}{r}(u_j)_r + \frac{1}{r^2}(u_j)_{\theta\theta} = 0.
\]
Conformality $(u_j)_z\cdot(u_j)_z=0$ translates to the orthogonality of coordinate lines:
\[
(u_j)_r\cdot(u_j)_\theta=0.
\]
The derivatives of each $u_j$ extend continuously up to the boundary arcs of $\hat{D}$ (the two radial junction segments $\gamma_1,\gamma_2$ and the free-boundary arc $\sigma$). Along each junction curve three sheets meet at $120^\circ$ angles (the local $Y$-balance condition), and the radial and second radial derivatives match across the junction:
\begin{equation}
\label{eq:matching-r}
u_i(p)=u_k(p)=u_l(p),\ \
(u_i)_r(p)=(u_k)_r(p)=(u_l)_r(p),\ \
(u_i)_{rr}(p)=(u_k)_{rr}(p)=(u_l)_{rr}(p)
\end{equation}
whenever $F_i,F_k,F_l$ meet along that curve. Substituting into the harmonicity equation yields
\begin{equation}
\label{eq:matching-theta2}
(u_i)_{\theta\theta}(p)=(u_k)_{\theta\theta}(p)=(u_l)_{\theta\theta}(p)
\end{equation}
along the junction. The $Y$-balance further gives
\begin{equation}
\label{eq:Y-balance}
(u_i)_\theta(p)+(u_k)_\theta(p)+(u_l)_\theta(p)=0.
\end{equation}
Along the free-boundary arc $\sigma$ we have $u_j(\sigma)\subset\mathbb{S}^{n-1}$ and the radial derivative is normal to the sphere, hence
\[
(u_j)_r=f\,u_j,
\]
for a scalar function $f$. Differentiating tangentially yields
\[
(u_j)_{r\theta}=f_\theta u_j + f\,(u_j)_\theta,
\]
so the normal component vanishes:
\begin{equation}
\label{eq:perp}
(u_j)_{r\theta}^\perp=0\qquad\text{along }\sigma.
\end{equation}
By the classical theory of conformal minimal immersions (see Theorem~1.1 in~\cite{matinpour2025cone}), the squared normal part of the second derivative
\[
Q_j^\perp(z):=(u_j)_{zz}^\perp
\]
satisfies $(Q_j^\perp)^2=(u_j)_{zz}^2$ and is therefore holomorphic on $\hat{D}$. In polar coordinates,
\[
u_{j,zz}=\frac14 e^{-2i\theta}\Bigl[(u_j)_{rr}-\frac{(u_j)_r}{r}-\frac1{r^2}(u_j)_{\theta\theta}-i\Bigl(\frac{2(u_j)_{r\theta}}{r}-\frac{2(u_j)_\theta}{r^2}\Bigr)\Bigr],
\]
and
\[
(u_j)_{zz}^\perp=\frac14 e^{-2i\theta}\Bigl[(u_j)_{rr}^\perp-\frac1{r^2}(u_j)_{\theta\theta}^\perp-\frac{2i}{r}(u_j)_{r\theta}^\perp\Bigr].
\]
Define the summed holomorphic quadratic differential on $\hat{D}$ by
\[
h(z)=\sum_{j=1}^6 Q_j^2(z)=\sum_{j=1}^6(Q_j^\perp)^2(z).
\]
A direct expansion gives
\[
h(z)=\frac{e^{-4i\theta}}{16}\sum_{j=1}^6\Bigl[(u_j)_{rr}-\frac{(u_j)_r}{r}-\frac1{r^2}(u_j)_{\theta\theta}-i\Bigl(\frac{2(u_j)_{r\theta}}{r}-\frac{2(u_j)_\theta}{r^2}\Bigr)\Bigr]^2
\]
and the analogous expression for the normal part. Set $H(z)=z^4 h(z)$. Its imaginary part reads
\[
\operatorname{Im} H(z)=\sum_{j=1}^6-\frac{ir^3}{4}(u_j)_{r\theta}^\perp\Bigl((u_j)_{rr}^\perp-\frac1{r^2}(u_j)_{\theta\theta}^\perp\Bigr).
\]
Condition~\eqref{eq:perp} implies that $H(z)$ is real along $\sigma$. Along the junction segments $\gamma_1\cup\gamma_2$ the matching conditions~\eqref{eq:matching-r}--\eqref{eq:matching-theta2} make the coefficient $(u_j)_{rr}-\frac{(u_j)_r}{r}-\frac1{r^2}(u_j)_{\theta\theta}$ independent of $j$, while the $Y$-balance~\eqref{eq:Y-balance} forces the sum of the coefficients $\frac{2(u_j)_{r\theta}}{r}-\frac{2(u_j)_\theta}{r^2}$ to vanish. Consequently $\operatorname{Im} H(z)=0$ along $\gamma_1\cup\gamma_2$ as well.\\

Thus $H(z)$ is holomorphic on $\hat{D}$ and real-valued on the entire boundary $\partial\hat{D}=\gamma_1\cup\gamma_2\cup\sigma$. By the Schwarz reflection principle, or equivalently, the fact that a holomorphic function real on the boundary of a disk is constant, $H(z)$ is constant. Since $H(0)=0$, we conclude $H\equiv0$ on $\hat{D}$. Inspecting the real and imaginary parts yields
\[
(u_j)_{r\theta}^\perp=0\qquad\text{and}\qquad(u_j)_{rr}^\perp-(u_j)_{\theta\theta}^\perp=0\qquad\text{along }\sigma.
\]
Minimality now implies $(u_j)_{rr}^\perp=(u_j)_{\theta\theta}^\perp=0$ along $\sigma$, so the second fundamental form of $\Sigma_j$ vanishes on the interior boundary arc $\partial'\Sigma_j:=\partial\Sigma_j\cap\mathbb{S}^{n-1}$. In particular the second fundamental form of $\partial'\Sigma_j$ in $\mathbb{S}^{n-1}$ vanishes, hence each $\partial'\Sigma_j$ is a geodesic arc (piece of a great circle) on the sphere.\\
Therefore the boundary $\partial\Sigma$ consists of geodesic arcs on $\mathbb{S}^{n-1}$ and forms a stationary geodesic network with tetrahedral combinatorics. Moreover $\Sigma$ is a free-boundary minimal $T$-surface spanning this network and meeting $\mathbb{S}^{n-1}$ orthogonally. Theorem~\ref{Thm:uniq. FBMT with NT} then yields the existence of an orthogonal transformation $A\in\mathrm{O}(n)$ such that $\Sigma=A(TC)$. This completes the proof of Theorem~\ref{thm:Tcone}.
\end{proof}

Theorem \ref{thm:uniq.PlateauSurfaces} then follows as a corollary to [Theorem 2.1 \cite{fraser2015uniqueness}], [Theorem 1.1 \cite{matinpour2025cone}], and  Theorem \ref{thm:Tcone}.

\begin{thebibliography}{15}

\bibitem{choe2025free}
J. ~Choe, 
\emph{Free Boundary Minimal Surfaces and the Reflection Principle},
arXiv preprint arXiv:2501.02700, (2025).

\bibitem{colding2011course}
T. H. ~Colding  and W. P. ~Minicozzi,
\emph{A Course in Minimal Surfaces}, volume \textbf{121}.
American Mathematical Soc., (2011).


\bibitem{fraser2015uniqueness}
A. ~Fraser and R. ~Schoen, 
\emph{Uniqueness Theorems for Free Boundary Minimal Disks in
Space Forms}, International Mathematics Research Notices, 2015(\textbf{17}):8268--8274, (2015).
Publisher: Oxford University Press.


\bibitem{nitsche1985stationary}
J. C. C. ~Nitsche, 
\emph{Stationary Partitioning of Convex Bodies},
Archive for rational mechanics and analysis, \textbf{89}:1--19, (1985).
Publisher: Springer.

\bibitem{matinpour2025cone}
E. ~Matinpour
 \emph{Rigidity and Index of Free Boundary Minimal Y-cones in the Unit Ball},
arXiv preprint
arXiv:2509.24137v2, (2025).




\end{thebibliography}


\end{document}

\section{Preliminaries}
Let $\Sigma_1,\ldots,\Sigma_6$ be compact smooth two–dimensional surfaces with boundary.
Let $\Gamma_1,\Gamma_2,\Gamma_3,\Gamma_4$ be smooth one–dimensional manifolds.\\
For each $j=1,\ldots,6$, suppose that the boundary of $\Sigma_j$ can be decomposed as
\[
\partial\Sigma_j=\partial'\Sigma_j\cup\partial''\Sigma_{j_1} \cup \partial'' \Sigma_{j_2} ,
\]
where $\partial'\Sigma_j,\ \partial''\Sigma_{j_1}$, and $\partial''\Sigma_{j_2}$ are disjoint smooth arcs
(possibly sharing some endpoints).\\
Assume that the arcs $\partial''\Sigma_{j_i}$ are identified in triples so as to form
four junction curves
\[
\Gamma_1,\Gamma_2,\Gamma_3,\Gamma_4 .
\]
More precisely, for each $n=1,\ldots,4$ there exist three faces
$\Sigma_{k},\Sigma_{l},\Sigma_{m}$ whose boundary arcs $\partial'' \Sigma_{k_i}, \partial''\Sigma_{l_i},\partial'' \Sigma_{m_i}$ are identified
via smooth diffeomorphisms with $\Gamma_n$, and more specifically, for each $j$ there are two $n_{j_1},n_{j_2}$ so that $\Gamma_{j_1}$ and $\Gamma_{j_2}$ are identified with $\partial'' \Sigma_{j_1}$ and $\partial'' \Sigma_{j_2}$, respectively. After these identifications, we obtain a surface
\[
\Sigma=\Big(\bigcup_{j=1}^6 \Sigma_j ; \Gamma_1,\Gamma_2,\Gamma_3,\Gamma_4\Big),
\]
Furthermore, assume that the four junction curves meet at a single point
\[
p\in \Sigma ,
\]
called the \emph{$T$--vertex}, so that the combinatorics of the six compact
faces and four junction curves agrees with the cone over the $1$–skeleton
of a regular tetrahedron.
Now, we call $\Sigma$ a \emph{$T$--singular surface} or simply a \emph{$T$--surface}.
The curves $\Gamma_k$ are the \emph{junction curves} along which three faces meet.\\
The boundary of $\Sigma$ is therefore
\[
\partial\Sigma=\bigcup_{j=1}^6 \partial'\Sigma_j .
\]


\begin{definition}
Let 
\[
\Sigma=\Big(\bigcup_{j=1}^6 \Sigma_j ; \Gamma_1,\Gamma_2,\Gamma_3,\Gamma_4\Big)
\]
be a $T$--surface.

A map
\[
\varphi=(\varphi_1,\ldots,\varphi_6):\Sigma\to\mathbb{R}^n
\]
is called a \emph{minimal immersion} if the following hold:

\begin{enumerate}

\item[(i)]
For each $j=1,\ldots,6$, the map
\[
\varphi_j:\Sigma_j\to\mathbb{R}^n
\]
is a smooth minimal immersion.

\item[(ii)]
Along each junction curve $\Gamma_k$ the three faces sharing that curve
have the same trace, producing a common immersed curve
\[
\gamma_k=\varphi(\Gamma_k).
\]

\item[(iii)]
Let $\tau_1,\tau_2,\tau_3$ denote the outward unit conormal vectors of the
three faces meeting along $\gamma_k$.
They satisfy the balance condition
\[
\tau_1+\tau_2+\tau_3=0 \quad \text{along } \gamma_k .
\]

\item[(iv)]
At the tetrahedral vertex $p$ the configuration of the six faces and four
junction curves is locally equivalent to the cone over the $1$–skeleton
of a regular tetrahedron centered at the origin.
\end{enumerate}

Such a surface is called a \emph{minimal $T$–surface}.
\end{definition}



\begin{definition}
The $T$--surface
\[
\Sigma=\Big(\bigcup_{j=1}^6 \Sigma_j ; \Gamma_1,\Gamma_2,\Gamma_3,\Gamma_4\Big)
\]
is called a \emph{free boundary minimal $T$–surface} in the unit ball
$\mathbb{B}^n$ if

\begin{enumerate}

\item each face $\Sigma_j$ is minimally immersed in $\mathbb{B}^n$,

\item the boundary
\[
\partial\Sigma=\bigcup_{j=1}^6\partial'\Sigma_j
\]
lies in the unit sphere $\mathbb{S}^{n-1}$,

\item the surfaces meet $\mathbb{S}^{n-1}$ orthogonally,

\item along each junction curve the conormals satisfy
\[
\tau_1+\tau_2+\tau_3=0,
\]

\item at the vertex the configuration is tetrahedral (a $T$–configuration).
\end{enumerate}
\end{definition}


\begin{definition}[Flat $T$–cone]\label{def:Tcone}
Let $P \subset \mathbb{R}^3$ be a regular tetrahedron centered at the origin,
with vertices $v_1,v_2,v_3,v_4 \in \mathbb{S}^2$.
Let $E(P)$ denote the $1$–skeleton of $P$, that is, the union of its six edges.

The \emph{flat $T$–cone} $TC$ is defined as the cone over $E(P)$ with vertex
at the origin:
\[
TC = \{ t x : x \in E(P),\ 0 \le t \le 1 \}.
\]

Equivalently, $TC$ is the union of six planar sectors,
each spanned by the origin and one edge of the tetrahedron.
The four line segments
\[
\Gamma_i = \{ t v_i : 0 \le t \le 1 \}, \qquad i=1,2,3,4,
\]
are called the \emph{junction curves}. Along each $\Gamma_i$ three faces meet
with equal angles of $120^\circ$.
\end{definition}
Let $TC = \bigcup_{j=1}^6 F_j$ denote the flat $T$–cone,
where $F_j$ are its six planar faces.\\
Parametrize each face $F_j$ by a quarter-disk $\hat D =\{ x^2+y^2 \leq 1,\ x\geq 0,\ y\geq 0 \}$. Set
\begin{equation}
\nonumber
\begin{split}
&\gamma_1=\{ (x,y)\in \mathbb{R}^2\ :\ x=0,\ 0\leq y\leq 1  \},\\
&\gamma_2=\{ (x,y)\in \mathbb{R}^2\ :\ y=0,\ 0\leq x\leq 1  \},\\
&\sigma=\{ (x,y)\in \mathbb{R}^2\ :\ x^2+y^2=1,\  x\geq 0,\ y\geq 0  \}.
\end{split}
\end{equation}
Suppose that $\partial''F_j=\Gamma_{j_1} \cup \Gamma_{j_2}$, then consider the parametrization above so that $\gamma_1,\ \gamma_2, \sigma$ parameterize $\Gamma_{j_1},\ \Gamma_{j_2}$ and $\partial' F_j$, respectively.
\medskip

Let $\Sigma = \big( \bigcup_{j=1}^6 \Sigma_j ; \Gamma_1,\Gamma_2,\Gamma_3,\Gamma_4 \big)$
be a $T$--surface with six faces and four junction curves.\\
The surface $\Sigma$ can be equipped with an abstract metric,
not necessarily induced by a Euclidean metric.
Let $g_j$ be a Riemannian metric on $\Sigma_j$.\\
We say that \emph{$(g_1,\ldots,g_6)$ defines a metric on $\Sigma$}
if the induced metrics agree along each junction curve.
More precisely, if three faces $\Sigma_{i},\Sigma_{j},\Sigma_{k}$
meet along $\Gamma_\ell$, then
\[
g_i|_{\Gamma_\ell} = g_j|_{\Gamma_\ell} = g_k|_{\Gamma_\ell}.
\]

Equivalently, all three faces induce the same metric
\[
g_{\Gamma_\ell}
\]
on the junction curve $\Gamma_\ell$.\\
Given such a metric $g$ on $\Sigma$, we define the unit outer conormal
vector field $\tau_j$ to $\Sigma_j$ along each junction curve.\\
Let $\eta$ be the unit tangent vector field along a junction curve
$\Gamma_\ell$.
If the three faces $\Sigma_i,\Sigma_j,\Sigma_k$ meet along $\Gamma_\ell$,
their geodesic curvatures are defined by
\[
\kappa_i = g_i(\nabla_{\eta}^{\Sigma_i} \eta ,\tau_i), \quad
\kappa_j = g_j(\nabla_{\eta}^{\Sigma_j} \eta ,\tau_j), \quad
\kappa_k = g_k(\nabla_{\eta}^{\Sigma_k} \eta ,\tau_k).
\]

\begin{definition}
We say that the metric $g$ on $\Sigma$ is \emph{compatible with the
$T$–structure} if along each junction curve $\Gamma_\ell$ the geodesic
curvatures satisfy
\[
\kappa_i + \kappa_j + \kappa_k = 0 .
\]
\end{definition}

If $\Sigma$ is a minimal $T$–surface in $\mathbb{R}^n$, then the sum of the
outward conormals of the three faces meeting along each junction curve
vanishes. Consequently the sum of the corresponding geodesic curvatures
also vanishes.


We equip $TC$ with the compatible metric $g=(g_1,\ldots,g_6)$
induced from the ambient Euclidean space.

\begin{definition}
A map
\[
u=(u_1,\ldots,u_6):TC \to \mathbb{R}^n
\]
is called \emph{conformal} if each
\[
u_j : F_j \to \mathbb{R}^n
\]
is conformal with respect to the induced metric and the maps agree
along the junction curves. That is,
if $p \in \Gamma_\ell$ and the faces $F_i,F_j,F_k$ meet along $\Gamma_\ell$,
then
\[
u_i(p)=u_j(p)=u_k(p).
\]
\end{definition}

\begin{definition}[Conformal map on a $T$--cone]
Let $TC=\bigcup_{j=1}^{6} F_j$ be the flat $T$--cone equipped with a
compatible metric
\[
g=(g_1,\dots,g_6),
\]
where each $g_j$ is a Riemannian metric on $F_j$ and the metrics agree
along the junction curves.

A map
\[
u=(u_1,\dots,u_6):TC\to\mathbb{R}^n
\]
is called \emph{conformal} if the following conditions hold.

\begin{enumerate}

\item For each face $F_j$, the restriction
\[
u_j : (F_j,g_j)\to\mathbb{R}^n
\]
is a conformal immersion; that is, there exists a positive function
$\lambda_j:F_j\to(0,\infty)$ such that
\[
(u_j)^*(g_{\mathbb{R}^n})=\lambda_j\, g_j ,
\]
where $g_{\mathbb{R}^n}$ denotes the Euclidean metric.

\item (Matching condition along junction curves)  
If three faces $F_i,F_j,F_k$ meet along a junction curve $\Gamma_\ell$,
then for every $p\in\Gamma_\ell$ we have
\[
u_i(p)=u_j(p)=u_k(p).
\]

\item (Compatibility of differentials along the junction)  
Along each junction curve $\Gamma_\ell$, the tangential derivatives
agree. If $\eta$ denotes the unit tangent vector field of $\Gamma_\ell$,
then
\[
du_i(\eta)=du_j(\eta)=du_k(\eta)
\quad \text{on } \Gamma_\ell .
\]

\end{enumerate}
\end{definition}

%%%%%%%%%%%%%%%%%%%%%%%%%%%%%%%%%%%%%%%%%%%%%%%%%%%%%%%%%%%%%%%%%%%%%%%%%%%%%%%%%%%%%%%%%%%%%%%%%%%%%%%%%%%%%%%%%%%%%%%%%%%%%%%%%%%%%%%%%%%%%%%%%%%%%%%%%%%%%%%%%%%%%%%%%%%%%%%%%%%%%%%%%%%%%%%%%%%%%%%%%%%%%%%%%%%%%%%%%%%%%%%%%%%%%%%%%%%%%%%%%%%%%%%%%%%%%%%%%%%%%%%%%%%%%%%%%%%%%%%%%%%%%%%%%%%%%%%%%%%%%%%%%%%%%%%%%%%%%%%%%%%%%%%%%%%%%%%%%%%%%%%%%%%%%%%%%%%%%%%%%%%%%%%%%%%%%%%%%%%%%%%%%%%%%%%%%%%%%%%%%%%%%%%%%%%%%%%%%%%%%%%%%%%%%%%%%%%%%%%%%%%%%%%%%%%%%%%%%%%%%%%%%%%%%%%%%%%%%%%%%%%%%%%%%%%%%%%%%%%%%%%%%%%%%%%%%%%%%%%%%%%%%%%%%%%%%%%%%%%%%%%%%%%%%%%%%%%%%%%%%%%%%%%%%%%%%%%%%%%%%%%%%%%%%%%%%%%%%%%%%%%%%%%%%%%%%%%%%%%%%%%%%%%


\begin{theorem}
Let \(\Sigma \subset B^n\) be a free-boundary minimal $2$-dimensional T-surface with exactly one T-point, four Y-junction edges (along each of which three faces meet at $120^\circ$), and six faces, each diffeomorphic to a quarter-disk. Assume that
\begin{itemize}
\item \(\partial \Sigma \cap \mathbb{S}^{n-1} \) is a stationary geodesic network with tetrahedral combinatorics,
\item \(\Sigma\) meets \(\mathbb{S}^{n-1}\) orthogonally.
\end{itemize}
Then there exists \(A \in \mathrm{O}(n)\) such that \(\Sigma = A(TC)\), where \(TC\) is the T-cone over the regular tetrahedron centered at the origin. In particular, the T-point is at the origin, the four Y-junctions are the straight radial segments from the origin to the vertices, and each face is a flat conical sector in the $2$-plane spanned by the corresponding pair of vertices.
\end{theorem}

\begin{proof}
Let \(N = \partial \Sigma \cap \mathbb{S}^{n-1}\) be the boundary network. By assumption \(N\) has tetrahedral combinatorics, each edge is a geodesic arc, and at each of its four vertices exactly three edges meet with the \(120^\circ\) equilibrium condition (the Y-configuration). Thus \(N\) is a stationary geodesic network with tetrahedral combinatorics. By the uniqueness lemma \ref{lem: uniq. of RT-network}, there exists \(A \in \mathrm{O}(n)\) such that \(N = A(\mathcal{RT})\), where \(\mathcal{RT}\) is the regular tetrahedral network on the great \(2\)-sphere in the subspace \(W \cong \mathbb{R}^3\) spanned by the standard regular-tetrahedron vertices \(v_1,\dots,v_4\) (with \(\langle v_i,v_j\rangle = -1/3\)). Without loss of generality (replacing \(\Sigma\) by \(A^{-1}(\Sigma)\)) we may assume that the boundary network is exactly the regular tetrahedral network \(\mathcal{RT} \subset W \cap \mathbb{S}^{n-1}\).
\medskip


Suppose $F$ is a face of $\Sigma$ with $\partial F=\sigma \cup \Gamma_1 \cup \Gamma_2$, where $\sigma$ is a geodesic arc on $\mathbb{RT}$ and $\Gamma_1,\Gamma_2$ are two junction curves of $\Sigma$. Let $\sigma\subset P$ be the plane passing through the origin. Then reflect $F$ and $P$ across the sphere along the boundary $\sigma$ via the inversion
$$
I(x)=\frac{x}{\vert x\vert^2}.
$$
Define $\widetilde{F}:=F\cup I(F).$
Then $\widetilde{F}$ is a smooth minimal surface in a neighborhood of $\sigma$, and $\sigma$ becomes an interior curve of $\widetilde{F}$. Since $P$ passes through the origin, we have $I(P)=P$. Thus, $P$ is also invariant under the reflection. By assumption, $\sigma \subset F \cap P$,
the outward conormal of $F$ agrees with that of $P$ along $\sigma$. This implies:
$$
T_pF=T_pP,\ \ \ \text{for all}\ p\in \sigma
$$
Thus, along $\sigma$, the surfaces $\widetilde{F}$ and $P$ intersect, and have the same tangent plane. So, they have second-order contact along $\sigma$. Therefore, by the unique continuation-strong maximum principle $\widetilde{F}\subset P$, and thus $F\subset P$.
\medskip

We know that there exists a regular $T$-cone $TC$ with $\partial TC=\mathbb{RT}$, and moreover, $TC$ has a flat face which passes through the origin and its boundary is $\sigma$. Since the argument holds for each face of $\Sigma$, and $\Sigma$ has the structure of a $T$-surface with a single $T$-point, four $Y$-junctions, and six faces, we conclude that $\Sigma=TC$. 
\medskip


Let \(p \in B^n\) be the unique T-point of \(\Sigma\). Near \(p\), \(\Sigma\) consists of six faces meeting along four Y-junction edges emanating from \(p\). Because \(\Sigma\) is minimal and the Y-junctions satisfy the \(120^\circ\) condition, the tangent cone of \(\Sigma\) at \(p\) is a stationary singular cone whose link (on the unit sphere in \(T_p\mathbb{R}^n\)) is a stationary geodesic network with tetrahedral combinatorics. Applying the uniqueness lemma once more to this link, the four outgoing unit tangent vectors \(\tau_1,\dots,\tau_4\) at \(p\) themselves form a regular tetrahedral configuration in some \(3\)-dimensional subspace \(V_p \subset T_p\mathbb{R}^n\), and the six local sectors are flat conical sectors in the six corresponding \(2\)-planes.
\medskip

The free-boundary condition (\(\Sigma\) meets \(\mathbb{S}^{n-1}\) orthogonally) together with the stationarity of the whole configuration implies that \(\Sigma\) cannot leave the \(3\)-dimensional subspace \(W\) containing the boundary: any component of a face transverse to \(W\) would produce a non-zero mean-curvature term or violate the balancing of conormals at the Y-junctions (by the maximum principle for minimal surfaces and the fact that the boundary data lie in \(W\)). Thus \(p \in W\) and \(V_p = W\) (the only \(3\)-dimensional subspace compatible with the boundary). Consequently the four outgoing directions \(\tau_k\) at \(p\) point precisely toward the four boundary vertices \(v_k\).
\medskip

Each Y-junction edge \(\gamma_k\) is therefore a curve from \(p\) to \(v_k\) whose initial direction is exactly the radial direction \(v_k - p\). By the triple-junction stationarity condition along \(\gamma_k\) (three minimal faces meeting at \(120^\circ\)) and the global minimality of \(\Sigma\), Lawlor's theorem (extended to free-boundary settings) forces each \(\gamma_k\) to be a straight line segment. Thus each \(\gamma_k\) is the radial segment \([p, v_k]\).
\medskip

It remains to show \(p = 0\). Suppose \(p \neq 0\). The free-boundary condition implies that the position vector field \(X(x) = x\) produces a vanishing first variation of area (the conormal along each boundary arc is parallel to the radial vector on \(\mathbb{S}^{n-1}\), and the junction balancing cancels the interior contributions). However, the monotonicity formula for free-boundary minimal surfaces in the ball shows that the area ratio \(r^{-2} \mathrm{Area}(\Sigma \cap B_r(0))\) is strictly increasing unless \(\Sigma\) is conical with vertex at \(0\). Since \(\Sigma\) has only one singularity (at \(p \neq 0\)) and its tangent cone at \(p\) is already conical, the only way the monotonicity ratio can be constant (forced by the first-variation vanishing for the radial field) is if \(p = 0\).
\medskip

With \(p = 0\) and the Y-junctions being the radial segments to the \(v_k\), each face is a minimal surface bounded by two radial segments and one great-circle arc on \(\mathbb{S}^{n-1}\). Because the face meets the sphere orthogonally and the two radial boundaries satisfy the \(120^\circ\) condition with the adjacent faces, the same holomorphic-quadratic-differential argument used for the Y-cone in Matinpour (2025) (applied sector by sector after conformal parametrization) shows that each face must be flat: the second fundamental form vanishes identically. Thus each face is precisely the flat conical sector in the plane spanned by the corresponding pair \(v_i,v_j\).

This is exactly the T-cone \(\mathcal{T}\). The proof is complete.
\end{proof}

%%%%%%%%%%%%%%%%%%%%%%%%%%%%%%%%%%%%%%%%%%%%%%%%%%%%%%%%%%%%%%%%%%%%%%%%%%%%%%%%%%%%%%%%%%%%%%%%%%%%%%%%%%%%%%%%%%%%%%%%%%%%%%%%%%%%%%%%%%%%%%%%%%%%%%%%%%%%%%%%%%%%%%%%%%%%%%%%%%%%%%%%%%

\begin{definition}[Conformal map on a $T$--cone]
Let $TC=\bigcup_{j=1}^{6} F_j$ be the flat $T$--cone equipped with a
compatible metric
\[
g=(g_1,\dots,g_6),
\]
where each $g_j$ is a Riemannian metric on $F_j$ and the metrics agree
along the junction curves.

A map
\[
u=(u_1,\dots,u_6):TC\to\mathbb{R}^n
\]
is called \emph{conformal} if the following conditions hold.

\begin{enumerate}

\item For each face $F_j$, the restriction
\[
u_j : (F_j,g_j)\to\mathbb{R}^n
\]
is a conformal immersion; that is, there exists a positive function
$\lambda_j:F_j\to(0,\infty)$ such that
\[
(u_j)^*(g_{\mathbb{R}^n})=\lambda_j\, g_j ,
\]
where $g_{\mathbb{R}^n}$ denotes the Euclidean metric.

\item (Matching condition along junction curves)  
If three faces $F_i,F_j,F_k$ meet along a junction curve $\Gamma_\ell$,
then for every $p\in\Gamma_\ell$ we have
\[
u_i(p)=u_j(p)=u_k(p).
\]

\item (Compatibility of differentials along the junction)  
Along each junction curve $\Gamma_\ell$, the tangential derivatives
agree. If $\eta$ denotes the unit tangent vector field of $\Gamma_\ell$,
then
\[
du_i(\eta)=du_j(\eta)=du_k(\eta)
\quad \text{on } \Gamma_\ell .
\]

\end{enumerate}
\end{definition}

%%%%%%%%%%%%%%%%%%%%%%%%%%%%%%%%%%%%%%%%%%%%%%%%%%%%%%%%%%%%%%%%%%%%%%%%%%%%%%%%%%%%%%%%%%%%%%%%%%%%%%%%%%%%%%%%%%%%%%%%%%%%%%%%
\begin{abstract}
    Nitsche--Fraser--Schoen~\cite{nitsche1985stationary},~\cite{fraser2015uniqueness} proved that any minimal disk satisfying the free boundary condition in the unit ball $\mathbb{B}^n$ must be an equatorial flat disk. The author \cite{matinpour2025cone} extended the uniqueness result to minimal $Y$-singular surfaces conformal to the flat $Y$-cone. In this paper, we further extend the uniqueness result to the minimal $T$-surfaces conformal to the flat $T$-cone meeting $\partial \mathbb{B}^n$ orthogonally. Therefore, we conclude that any conformal and minimal immersion of a flat model cone meeting $\partial \mathbb{B}^n$ orthogonally must itself be the flat model cone, which establishes a uniqueness result on free boundary minimal Plateau surfaces under conformality assumption with the reference model cone.
\end{abstract}

%%%%%%%%%%%%%%%%%%%%%%%%%%%%%%%%%%%%%%%%%%%%%%%%%%%%%%%%%%%%%%%%%%%%%%%%%%%%%%%%%%%%%%%%%%%%%%%%%%%%%%%%
Conformal Free-Boundary Minimal Model Cones in the Unit Ball are Rigid---Rigidity of Model Cones among Conformal Free-Boundary Minimal Surfaces in the Unit Ball----Free-Boundary Minimal Model Cones are Rigid under Conformal Immersions
\\


%%%%%%%%%%%%%%%%%%%%%%%%%%%%%%%%%%%%%%%%%%%%%%%%%%%%%%%%%%%%%%%%%%%%%%%%%%%%%%%%%%%%%%%%%%%%%%%%%%%%%%%%%%%%%%%%%%%%%%%%%%%%%%%%%%%%%%%%%%%%%%%%%%%%%%%%%%%%%%%%%%%%%%%%%%%%%%

\begin{proof}
The analysis of the Hopf differential is identical in spirit to that of the $Y$-cone case in~\cite{matinpour2025cone}. 
We therefore focus on the aspects that are specific to the tetrahedral geometry.\\


Define the normal component of the second derivative
\[
{u_j}_{zz}^{\perp} = \text{component of } {u_j}_{zz} \text{ normal to } \Sigma_j.
\]
Then it follows from the classical identity for conformal minimal immersions that $({u_j}_{zz}^{\perp})^2=({u_j}_{zz})^2$ and thus $({u_j}_{zz}^{\perp})^2$ is holomorphic on $\hat D$, see  [Theorem 1.1, \cite{matinpour2025cone}].\\
In the polar coordinates, on $\hat{D}$:
 $$
 {u_j}_{zz} =\frac{1}{4} e^{-2i\theta}\Big[  {u_j}_{rr}-\frac{{u_j}_r}{r} - \frac{1}{r^2} {u_j}_{\theta \theta} - i(\frac{2{u_j}_{r\theta}}{r}-\frac{2{u_j}_{\theta}}{r^2}) \Big],
 $$
 and,
 $$
{u_j}_{zz}^{\perp} =\frac{1}{4} e^{-2i\theta}\Big[  {u_j}_{rr}^{\perp} - \frac{1}{r^2} {u_j}_{\theta \theta}^{\perp} - \frac{2i}{r} {u_j}_{r\theta}^{\perp} \Big].
$$
Suppose that $F_{i},F_{k},F_{l}$ are joined along $\gamma_1$ (resp. $\gamma_2$), then on $\gamma_1$ (resp. $\gamma_2$)
$$
u_{i}(p)=u_{k}(p)=u_{l}(p),
$$
and,
\begin{equation} \label{eqns: p.derivs(r,rr)}
    \begin{split}
       &  {u_i}_{r}(p) = {u_k}_{r}(p) = {u_l}_{r}(p),  \\
       & {u_i}_{rr}(p) = {u_k}_{rr}(p) = {u_l}_{rr}(p).
    \end{split}
\end{equation}
Using equations (\ref{eqns: p.derivs(r,rr)}) in $\Delta u_j ={u_j}_{rr}+\frac{1}{r}{u_j}_r+\frac{1}{r^2}{u_j}_{\theta \theta}=0$ we find that
\begin{equation} \label{eqn: p.deriv(theta,theta)}
  {u_i}_{\theta \theta}(p) ={u_k}_{\theta \theta}(p) ={u_l}_{\theta \theta} (p),  
\end{equation}
along $\gamma_1$.
Moreover, by the $Y$-configuration along $\gamma_1$ we also have
\begin{equation}\label{eqn:p.deriv(theta)}
 {u_i}_{\theta}(p) +  {u_k}_{\theta}(p) + {u_l}_{\theta}(p) =0.
\end{equation}
 Set $Q_j(z)={u_j}_{zz}(z)$, and note that ${Q_j^{\perp}}^2(z)=Q_j^2(z)$. Define
$$
h(z)=\underset{j=1}{\overset{6}{\sum}}Q_j^2(z)=\underset{j=1}{\overset{6}{\sum}}{Q_j^{\perp}}^2(z),
$$
on $\hat{D}$. Compute
\begin{equation}
\nonumber
\begin{split}
h(z)&=\underset{j=1}{\overset{6}{\sum}}Q_j^2(z)=\underset{j=1}{\overset{6}{\sum}} \Big ( \frac{1}{16}e^{-4i\theta}   \big [{u_j}_{rr}-\frac{{u_j}_r}{r} -\frac{1}{r^2}{u_j}_{\theta \theta}-i(\frac{2{u_j}_{r\theta}}{r}-\frac{{2u_j}_{\theta}}{r^2}) \big ]^2 \Big )\\
&= \frac{e^{-4i\theta}}{16} \underset{j=1}{\overset{6}{\sum}} \Big( ({u_j}_{rr}-\frac{{u_j}_r}{r} -\frac{1}{r^2}{u_j}_{\theta \theta} )^2 -(\frac{2{u_j}_{r\theta}}{r}-\frac{2{u_j}_{\theta}}{r^2})^2\\
&\ -2i (\frac{2{u_j}_{r\theta}}{r}-\frac{2{u_j}_{\theta}}{r^2}) ({u_j}_{rr}-\frac{{u_j}_r}{r}-\frac{1}{r^2}{u_j}_{\theta \theta} ) \Big),
\end{split}
\end{equation}
and similarly,
\begin{equation}
    \nonumber
\begin{split}
h(z)& =\underset{j=1}{\overset{6}{\sum}}{Q_j^{\perp}}^2(z)=  \frac{1}{16} e^{-4i\theta} \underset{j=1}{\overset{6}{\sum}}\Big[  {u_j}_{rr}^{\perp} - \frac{1}{r^2} {u_j}_{\theta \theta}^{\perp} - \frac{2i}{r} {u_j}_{r\theta}^{\perp} \Big]^2 \\ 
&=\frac{e^{-4i\theta}}{16} \underset{j=1}{\overset{6}{\sum}} \Big( ({u_j}_{rr}^{\perp}-\frac{1}{r^2}{u_j}_{\theta \theta}^{\perp} )^2 -\frac{4}{r^2}{{u_j}_{r\theta}^{\perp}}^2-\frac{4i}{r} {u_j}_{r\theta}^{\perp} ({u_j}_{rr}^{\perp}-\frac{1}{r^2}{u_j}_{\theta \theta}^{\perp} ) \Big).
\end{split}
\end{equation}
Set $H(z)=z^4h(z)$. The imaginary part of $H(z)$ is
\begin{equation}
    \nonumber
    \begin{split}
        \operatorname{Im} H(z)&=\underset{j=1}{\overset{6}{\sum}}-\frac{ir^4}{8} (\frac{2{u_j}_{r\theta}}{r}-\frac{2{u_j}_{\theta}}{r^2}) ( {u_j}_{rr}-\frac{{u_j}_r}{r} -\frac{1}{r^2}{u_j}_{\theta \theta} )\\
        &=\underset{j=1}{\overset{6}{\sum}}-\frac{ir^3}{4} {u_j}_{r\theta}^{\perp} ({u_j}_{rr}^{\perp}-\frac{1}{r^2}{u_j}_{\theta \theta}^{\perp} ).
 \end{split}
\end{equation}
Since ${u_j}_{r\theta}^{\perp}=0$ on $\sigma$, see (\ref{eq:perp}), $H(z)$ is real on $\sigma$. \\
On the other hand, it follows from equations (\ref{eqns: p.derivs(r,rr)}) and (\ref{eqn: p.deriv(theta,theta)}) that the terms ${u_j}_{rr}-\frac{{u_j}_r}{r} -\frac{1}{r^2}{u_j}_{\theta \theta}$ in $\operatorname{Im} H(z)$ are independent of $j$ along $\gamma_1\cup \gamma_2$. By the $Y$-configuration, see (\ref{eqn:p.deriv(theta)}), the terms $(\frac{2{u_j}_{r\theta}}{r}-\frac{2{u_j}_{\theta}}{r^2})$ sum up to zero along $\gamma_1\cup \gamma_2$, and thus
 $\operatorname{Im} H(z)$ is equal to zero along $\gamma_1\cup \gamma_2$.\\
Therefore, we have that $H(z)$ is holomorphic on $\hat{D}$ and is real on $\partial \hat{D}=\gamma_1\cup \gamma_2\cup \sigma$. Hence, it must be constant on $\hat{D}$, and since it vanishes at zero, it must vanish on $\hat{D}$. By checking on the real and imaginary parts of $H(z)$, this implies that ${u_j}_{r\theta}^{\perp}=0$ and  ${u_j}_{rr}^{\perp}-{u_j}_{\theta \theta}^{\perp}=0$ on $\sigma$.\\
By the minimality, this implies that ${u_j}_{rr}^{\perp}={u_j}_{\theta \theta}^{\perp}=0$ on $\sigma$. Therefore, the second fundamental form of $\Sigma_j$ is zero on $\partial' \Sigma_j$. It follows that the second fundamental form of $\partial' \Sigma_j$ in $\mathbb{S}^{n-1}$ is zero. Hence, $\partial' \Sigma_j$ is a piece of a great circle, and thus $\partial \Sigma$ is composed of pieces of great circles. Therefore, by Theorem \ref{Thm:uniq. FBMT with NT}, $\Sigma =A(TC)$, for some $A\in O(n)$, and the proof is complete.
\end{proof}

%%%%%%%%%%%%%%%%%%%%%%%%%%%%%%%%%%%%%%%%%%%%%%%%%%%%%%%%%%%%%%%%%%%%%%%%%%%%%%%%%%%%%%%%%%%%%%%%%%%%%%%%%%%%%%%%%%%%%%%%%%%%%%%%%%%%%%%%%%%

\textit{Fourth,} a direct application of the spherical law of cosines yields 
$\cos \sigma = -\frac{1}{3}$, so that $\sigma = \arccos(-1/3)$.\\
we show that four equidistant points on $\mathbb{S}^2$ form a regular tetrahedron.\\ 
It is a classical result that, up to isometry of $\mathbb{S}^2$, there is only one configuration of four points on the sphere such that every pair is at the same distance $\sigma \in (0,\pi)$: the set of vertices of a regular tetrahedron inscribed in $\mathbb{S}^2$. To determine the precise value of $\sigma$, apply the spherical law of cosines in any of the equilateral triangles with side $\sigma$ and angle $120^\circ$:
\[
\cos \sigma = \cos \sigma \cdot \cos \sigma + \sin \sigma \cdot \sin \sigma \cdot \cos 120^\circ = \cos^2 \sigma + \sin^2 \sigma \cdot \Bigl(-\frac{1}{2}\Bigr).
\]
Rearranging gives
\[
\cos \sigma = \cos^2 \sigma - \frac{1}{2} (1 - \cos^2 \sigma) = \frac{3}{2} \cos^2 \sigma - \frac{1}{2}.
\]
Let $x = \cos \sigma$. Then
\[
3x^2 - 2x - 1 = 0 \quad \Longrightarrow \quad x = \frac{2 \pm \sqrt{4 + 12}}{6} = \frac{2 \pm 4}{6}.
\]
The roots are $x = 1$ (corresponding to the degenerate case $\sigma = 0$) and
\[
x = -\frac{1}{3} \quad \Longrightarrow \quad \sigma = \arccos\Bigl(-\frac{1}{3}\Bigr).
\]
The value $\arccos(-1/3)$ is precisely the spherical distance between any two vertices of a regular tetrahedron inscribed in the unit sphere, for example, with vertices at suitably normalized vectors $(1,1,1)$, $(1,-1,-1)$, $(-1,1,-1)$, $(-1,-1,1)$ whose pairwise inner products equal $-1/3$.
\medskip

\textit{Fourth,} we determine the value of the common side length and identify the configuration. \\ 
It is a classical fact that, up to isometry of $\mathbb{S}^2$, there is a unique configuration of four points such that every pair is at the same spherical distance $\sigma \in (0,\pi)$, namely the vertices of a regular tetrahedron. 

To compute $\sigma$, consider one of the equilateral spherical triangles with side length $\sigma$ and angles $120^\circ$. By the spherical law of cosines,
\[
\cos \sigma = \cos^2 \sigma + \sin^2 \sigma \cos(120^\circ).
\]
Since $\cos(120^\circ) = -\tfrac{1}{2}$, this simplifies to
\[
\cos \sigma = \tfrac{3}{2}\cos^2 \sigma - \tfrac{1}{2},
\]
which yields $\cos \sigma = -\tfrac{1}{3}$ (excluding the degenerate solution $\cos \sigma = 1$). Hence
\[
\sigma = \arccos\Bigl(-\tfrac{1}{3}\Bigr).
\]

Therefore the four vertices form a regular tetrahedron, and all edges of $\mathcal{N}$ have the same spherical length $\arccos(-1/3)$.
\medskip

\textit{Third,} we show that each spherical triangle is equilateral. \\ 
Consider any one of the four spherical triangles with angles $A = B = C = 120^\circ = 2\pi/3$. By the spherical law of sines,
\[
\frac{\sin a}{\sin A} = \frac{\sin b}{\sin B} = \frac{\sin c}{\sin C}.
\]
Since $A = B = C$, we have $\sin a = \sin b = \sin c$. Because the edges of $\mathcal{N}$ are proper geodesic arcs between distinct vertices and form a simple embedding, that means there is no self-intersections except at vertices, each edge length belongs to the open interval $(0, \pi)$. In this range the function $\sin$ is one-to-one on $(0,\pi)$, so
\[
\sin a = \sin b = \sin c \quad \Longrightarrow \quad a = b = c.
\]
Thus every spherical triangle is equilateral. Let $\sigma \in (0,\pi)$ be the common spherical side length of all four triangles. Since every edge of $\mathcal{N}$ is a side of at least one triangle, all six edges of $\mathcal{N}$ have the same length $\sigma$. In other words, every pair of vertices of $\mathcal{N}$ is at spherical distance exactly $\sigma$.
\medskip
